\documentclass[10pt,a4paper]{amsart}
\usepackage[margin=19mm]{geometry}
\usepackage{amsmath,amssymb,mathtools}
\usepackage[scr=rsfs,cal=euler]{mathalfa}
\usepackage{xfrac}
\usepackage[unicode,hidelinks,bookmarks=false]{hyperref}
\newcommand{\ie}{i.e.\ }
\newcommand{\cf}{cf.\ }
\newcommand{\resp}{respectively\ }
\newcommand{\Subsection}{\subsection}
\providecommand{\nobreakdash}{\nobreak\hbox{-}}
\setlength{\emergencystretch}{2em}
\multlinegap0pt
\usepackage{bm}
\usepackage{bbm}
\usepackage{dsfont}
\usepackage{xspace}
\usepackage{cancel}
\usepackage{mathtools}
\usepackage{amsthm}
\makeatletter
\@ifundefined{fact}{\newtheorem{fact}{Fact}}{}
\@ifundefined{lemma}{\newtheorem{lemma}[fact]{Lemma}}{}
\makeatother
\title[Higher order double point formulas via SSM-Thom polynomials]{Higher order double point formulas via SSM-Thom polynomials}
\author{Reese Lance}
\date{Source: arXiv:2601.17651v1 (CC BY 4.0)}
\begin{document}
\maketitle

\noindent\emph{Source and licence.} Higher order double point formulas via SSM-Thom polynomials. Version arXiv:2601.17651v1 (CC BY 4.0), distributed under the Creative Commons Attribution 4.0 International licence, as recorded in the arXiv licence metadata for this exact version and shown on the abstract page https://arxiv.org/abs/2601.17651v1. Full rights notice: licenses/arxiv-2601.17651-CC-BY-4.0.txt.\par\smallskip

\noindent\textit{Editorial scope.} Counted unit: the Lemma whose source label is \texttt{Sm determinant} in the article (source0.tex lines 1180--1184, source label \texttt{Sm determinant}), delivered together with the article's own complete original proof of it (source0.tex lines 1185--1325, one \texttt{proof} environment of 141 lines). No mathematics is written, repaired or re-derived: statement and proof are the article's own text line for line. Required context retained uncounted: none. Every definition, display and label of those lines is retained unchanged. Omitted from this package and not redistributed: 1--1179, 1326--1465. Nothing inside a retained excerpt is omitted or altered: every retained body is delivered exactly as the article's lines are written, so no omission receipt of any kind accompanies this package. Citations: the retained excerpts carry no citation command at all, so no bibliography entry of the article is transcribed and no reference is redistributed. Reference adaptation: none was needed and none was made; every reference inside the delivered unit resolves inside it, so no unresolved reference or citation is produced. Printed slips kept literal: no printed word, symbol, hypothesis or number of the counted statement or of its proof was altered. Layout and class: the article's own class is \texttt{article} with packages including inputenc, babel, amsmath, amsfonts, amssymb, amsthm, tikz-cd, mathtools, geometry, enumerate, enumitem, mathrsfs, hyperref, graphicx; the wrapper is the lane's pinned \texttt{amsart} editorial wrapper at 10pt on A4, which restates the theorem-like environments the retained text needs and adds the article's own macro definitions that the retained text needs (none), with extra wrapper packages (bm, bbm, dsfont, xspace, cancel, mathtools). The delivered numbering is the unit's own. Assets: no retained span contains a figure environment, a graphics-inclusion command, a PDF-page inclusion, a table or a TikZ picture; no image file is copied into this package, the package declares no asset dependency and no image file is redistributed with it. Licence: version arXiv:2601.17651v1 of this article is distributed under the Creative Commons Attribution 4.0 International licence, as recorded in the arXiv licence metadata for this exact version and shown on the abstract page https://arxiv.org/abs/2601.17651v1. A full rights notice is delivered at licenses/arxiv-2601.17651-CC-BY-4.0.txt. Characters: every retained excerpt is pure ASCII, so the delivered engine, XeLaTeX with its default Latin Modern text font, reports no missing character.
\begin{lemma}\label{Sm determinant}
\[
S_{(m)}=\left[\sum_{i=1}^{m}2^{i-1}\cdot \alpha^i\cdot h_{m-i}(\alpha+\beta_1,\dots,\alpha+\beta_\ell)\right] + h_{m}(\alpha+\beta_1,\dots,\alpha+\beta_\ell).
\]
\end{lemma}

\begin{proof}
Using the cofactor expansion of determinant once using the first row, we obtain 
\begin{gather*}
S_{(m)}(c_1,\dots,c_m) \equiv \begin{vmatrix}c_1 & c_2 & c_3&\dots&c_{m-1} & c_{m}\\ 1 & c_1 &c_2& \dots&c_{m-2} & c_{m-1}\\0&1&c_1&\dots\\ & \vdots & \\
&&&&1&c_1\end{vmatrix}\\
=c_1\begin{vmatrix}c_1 & c_2 &c_3& \dots&c_{m-2} & c_{m-1}\\ 1 & c_1 & c_2&\dots&c_{m-3} & c_{m-2}\\0&1&c_1&\dots\\ & \vdots & \\
&&&&1&c_1\end{vmatrix}-c_2\begin{vmatrix}1 & c_2 & c_3 &\dots&c_{m-2} & c_{m-1}\\ 0& c_1 & c_2 & \dots&c_{m-2} & c_{m-3}\\0&1&c_1&\dots\\ & \vdots & \\
&&&&1&c_1\end{vmatrix} \\
+ \dots + (-1)^{m+1}c_m\begin{vmatrix}1 & c_1 & c_2 &\dots&c_{m-3} & c_{m-2}\\ 0& 1 & c_1 & \dots&c_{m-4} & c_{m-3}\\0&0&1&\dots\\ & \vdots & \\
&&&&0&1\end{vmatrix}.
\end{gather*}
The first term is exactly $c_1S_{(m-1)}$. The second may be computed by again performing a cofactor expansion, this time on the first column, since there is only one non-zero entry. The second term then becomes $-c_2\cdot S_{(m-2)}$. The $n$th term in the sum yields $(-1)^{n+1}c_nS_{(m-n)}$ after performing $n-1$ cofactor expansions. We obtain
$$
S_{(m)} = \sum_{i=1}^m(-1)^{i+1} c_i S_{(m-i)}.
$$
By strong induction we may assume 
$$
S_{(m-i)} = \left[\sum_{k=1}^{(m-i)}2^{k-1}\cdot \alpha^k \cdot h_{(m-i)-k}(\alpha+\beta_1,\dots,\alpha+\beta_\ell)\right]+h_{(m-i)}(\alpha+\beta_1,\dots,\alpha+\beta_\ell)
$$
for all $i \in \{1,2,\dots,m\}$, thus (For now, all symmetric polynomials have arguments $\alpha+\beta_1,\dots,\alpha+\beta_\ell$, so we abbreviate to simply $e_{(-)}$ and $h_{(-)}$)
\begin{gather*}
S_{(m)} = \sum_{i=1}^m(-1)^{i+1}c_i\left( \left[\sum_{k=1}^{(m-i)}2^{k-1}\cdot \alpha^k \cdot h_{(m-i)-k}\right]+h_{(m-i)}\right)\\
= \sum_{i=1}^m(-1)^{i+1}\left(e_i+(-1)^{i+1}\cdot \alpha^i + \sum_{a=1}^{i-1}e_a\cdot (-1)^{-(i-a)+1}\alpha^{i-a}\right)\cdot\left( \left[\sum_{k=1}^{(m-i)}2^{k-1}\cdot \alpha^k \cdot h_{(m-i)-k}\right]+h_{(m-i)}\right)\\\\
=\sum_{i=1}^m(-1)^{i+1}\Bigg(e_i \cdot \sum_{k=1}^{m-i}2^{k-1}\alpha^kh_{m-i-k} + e_i h_{m-i} + (-1)^{i+1}\alpha^i\cdot \sum_{k=1}^{m-i}2^{k-1}\alpha^kh_{m-i-k} + (-1)^{i+1}\alpha^i h_{m-i} \\
+ \sum_{a=1}^{i-1}e_a (-1)^{-(i-a)+1}\alpha^{i-a}\cdot \sum_{k=1}^{m-i}2^{k-1}\alpha^kh_{m-i-k} + \sum_{a=1}^{i-1}e_a(-1)^{-(i-a)+1}a^{i-a}h_{m-i}\Bigg).
\end{gather*}
Pulling the second term out of the parenthesis, then out of the sum, we get
\begin{gather*}
= \sum_{i=1}^m(-1)^{i+1} e_ih_{m-i} + \sum_{i=1}^m(-1)^{i+1}\Bigg(e_i \cdot \sum_{k=1}^{m-i}2^{k-1}\alpha^kh_{m-i-k} + (-1)^{i+1}\alpha^i\cdot \sum_{k=1}^{m-i}2^{k-1}\alpha^kh_{m-i-k} + (-1)^{i+1}\alpha^i h_{m-i} \\
+ \sum_{a=1}^{i-1}e_a (-1)^{-(i-a)+1}\alpha^{i-a}\cdot \sum_{k=1}^{m-i}2^{k-1}\alpha^kh_{m-i-k} + \sum_{a=1}^{i-1}e_a(-1)^{-(i-a)+1}\alpha^{i-a}h_{m-i}\Bigg).
\end{gather*}
The identity 
$$
\sum_{i=0}^m(-1)^ie_i h_{m-i}=0
$$
implies that the first sum above may be replaced with $h_m$, since the sum is missing the $i=0$ term:
\begin{gather*}
= h_m+ \sum_{i=1}^m(-1)^{i+1}\Bigg(e_i \cdot \sum_{k=1}^{m-i}2^{k-1}\alpha^kh_{m-i-k} + (-1)^{i+1}\alpha^i\cdot \sum_{k=1}^{m-i}2^{k-1}\alpha^kh_{m-i-k} + (-1)^{i+1}\alpha^i h_{m-i} \\
+ \sum_{a=1}^{i-1}e_a (-1)^{-(i-a)+1}\alpha^{i-a}\cdot \sum_{k=1}^{m-i}2^{k-1}\alpha^kh_{m-i-k} + \sum_{a=1}^{i-1}e_a(-1)^{-(i-a)+1}\alpha^{i-a}h_{m-i}\Bigg).
\end{gather*}
Splitting the sums into terms involving $e$'s and terms not involving $e$'s, we obtain
\begin{gather*}
= h_m+ \underline{\sum_{i=1}^m(-1)^{i+1}\Bigg(e_i \cdot \sum_{k=1}^{m-i}2^{k-1}\alpha^kh_{m-i-k} + \sum_{a=1}^{i-1}e_a (-1)^{-(i-a)+1}\alpha^{i-a}\cdot \sum_{k=1}^{m-i}2^{k-1}\alpha^kh_{m-i-k}} \\
\underline{+h_{m-i}\sum_{a=1}^{i-1}e_a(-1)^{-(i-a)+1}\alpha^{i-a}\Bigg)}\\
+\sum_{i=1}^m\left(\left[\sum_{k=1}^{m-i}2^{k-1}\alpha^{k+i}h_{m-i-k}\right]+\alpha^ih_{m-i}\right)
\end{gather*}
Now we claim that the underlined portion is equal to 0:

\begin{lemma}\label{sum=0}
$$
\sum_{i=1}^m(-1)^{i+1}\Bigg(e_i \cdot \sum_{k=1}^{m-i}2^{k-1}\alpha^kh_{m-i-k} + \sum_{a=1}^{i-1}e_a (-1)^{-(i-a)+1}\alpha^{i-a}\cdot \sum_{k=1}^{m-i}2^{k-1}\alpha^kh_{m-i-k}
$$
$$
+h_{m-i}\sum_{a=1}^{i-1}e_a(-1)^{-(i-a)+1}\alpha^{i-a}\Bigg)=0
$$
\end{lemma}

\begin{proof}
We proceed by induction on $m$: Let $P(m)$ be the proposition that the above sum is equal to 0. Checking the base case is routine, but there is one thing to note: That the sum is equal to 0 can be obtained \textit{as a polynomial in the $\alpha$, $e$'s and $h$'s. In particular, we do not need to refer to any properties of the $e$'s and $h$'s: we do not even need to assume that $h_0=1$. This implies that if we rename all the variables, the identity will still hold. This will be relevant later.}

We must show that 
$$
\sum_{i=1}^{m+1}(-1)^{i+1}\Bigg(e_i \cdot \sum_{k=1}^{m+1-i}2^{k-1}\alpha^kh_{m+1-i-k} + \sum_{a=1}^{i-1}e_a (-1)^{-(i-a)+1}\alpha^{i-a}\cdot \sum_{k=1}^{m+1-i}2^{k-1}\alpha^kh_{m+1-i-k}
$$
$$
+h_{m+1-i}\sum_{a=1}^{i-1}e_a(-1)^{-(i-a)+1}\alpha^{i-a}\Bigg)=0.
$$
Separating out the $i=m+1$ term (many of the inner sums are thus 0 because they start at $k=1$ and $a=1$), we get
\begin{gather*}
=\underbrace{(-1)^m\cdot h_0 \sum_{a=1}^me_a\cdot(-1)^{a-m}\cdot \alpha^{m+1-a}}_{i=m+1}\\
+\sum_{i=1}^{m}(-1)^{i+1}\Bigg(e_i \cdot \sum_{k=1}^{m+1-i}2^{k-1}\alpha^kh_{m+1-i-k} + \sum_{a=1}^{i-1}e_a (-1)^{-(i-a)+1}\alpha^{i-a}\cdot \sum_{k=1}^{m+1-i}2^{k-1}\alpha^kh_{m+1-i-k}\\
+h_{m+1-i}\sum_{a=1}^{i-1}e_a(-1)^{-(i-a)+1}\alpha^{i-a}\Bigg).
\end{gather*}
Then we separate out the final term of the inner sums that involve $m+1$
\begin{gather*}
= h_0 \sum_{a=1}^me_a\cdot(-1)^a\cdot \alpha^{m+1-a}\\
+\sum_{i=1}^{m}(-1)^{i+1}\Bigg(e_i \cdot \left[\sum_{k=1}^{m-i}2^{k-1}\alpha^kh_{m+1-i-k} +\underbrace{2^{m-i}\alpha^{m+1-i}h_0}_{k=m+1-i}\right]\\
+ \sum_{a=1}^{i-1}e_a (-1)^{-(i-a)+1}\alpha^{i-a}\cdot \left[\sum_{k=1}^{m-i}2^{k-1}\alpha^kh_{m+1-i-k}+\underbrace{2^{m-i}\alpha^{m+1-i}h_0}_{k=m+1-i}\right]\\
+h_{m+1-i}\sum_{a=1}^{i-1}e_a(-1)^{-(i-a)+1}\alpha^{i-a}\Bigg).
\end{gather*}
Then pull these extra $k=m+1-i$ terms all the way out of the outermost sum, upon which they pick up a sum over $i$ and an alternating sign,
\begin{gather*}
=h_0\sum_{a=1}^me_a\cdot(-1)^a\cdot \alpha^{m+1-a}\\
+\sum_{i=1}^m(-1)^{i+1}\left(e_i\cdot 2^{m-i}\alpha^{m+1-i}h_0 + \sum_{a=1}^{i-1}e_a(-1)^{-(i-a)+1}\alpha^{i-a}\cdot(2^{m-i}\alpha^{m+1-i}h_0)\right)\\
+\sum_{i=1}^{m}(-1)^{i+1}\Bigg(e_i \cdot \left[\sum_{k=1}^{m-i}2^{k-1}\alpha^kh_{m+1-i-k}\right]\\
+ \sum_{a=1}^{i-1}e_a (-1)^{-(i-a)+1}\alpha^{i-a}\cdot \left[\sum_{k=1}^{m-i}2^{k-1}\alpha^kh_{m+1-i-k}\right]\\
+h_{m+1-i}\sum_{a=1}^{i-1}e_a(-1)^{-(i-a)+1}\alpha^{i-a}\Bigg).
\end{gather*}
We want to cancel the final 3 lines in the expression above using the inductive hypothesis. However the sum in $P(m)$ involves variables $h_0,\dots,h_{m-i-1}$, while the expression above has variables $h_1,\dots,h_{m-i}$, due to the $+1$ in all $h$ subscripts. As argued before, $P(m)$ is true for any indeterminate variables, so we can still apply the inductive hypothesis to conclude the sum of the final 3 lines are 0,
\begin{gather*}
=h_0 \sum_{a=1}^me_a\cdot(-1)^a\cdot \alpha^{m+1-a}\\
+\sum_{i=1}^m(-1)^{i+1}2^{m-i}\alpha^{m+1-i}h_0\left(e_i+\sum_{a=1}^{i-1}e_a(-1)^{-(i-a)+1}\alpha^{i-a}\right)\\\\
=h_0 \sum_{a=1}^me_a\cdot(-1)^a\cdot \alpha^{m+1-a}\\
+\sum_{i=1}^m(-1)^{i+1}2^{m-i}\alpha^{m+1-i}h_0e_i\\
+\sum_{i=1}^m(-1)^{i+1}2^{m-i}\alpha^{m+1-i}h_0\sum_{a=1}^{i-1}e_a(-1)^{-(i-a)+1}\alpha^{i-a}.
\end{gather*}
Every term has a factor of $h_0$. Since we seek to show that this is equal to 0, we can factor and ignore the $h_0$ (note again that we do not need to impose that $h_0=1)$
\begin{gather*}
=\sum_{a=1}^me_a\cdot(-1)^a\cdot \alpha^{m+1-a}+\sum_{i=1}^m(-1)^{i+1}2^{m-i}\alpha^{m+1-i}e_i+\sum_{i=1}^m\sum_{a=1}^{i-1}e_a(-1)^a2^{m-i}\alpha^{m+1-a}.
\end{gather*}
In the third term, we swap the order of summation
\begin{gather*}
= \sum_{a=1}^me_a\cdot(-1)^a\cdot \alpha^{m+1-a}+\sum_{i=1}^m(-1)^{i+1}2^{m-i}\alpha^{m+1-i}e_i+\sum_{a=1}^{m-1}\sum_{i=a+1}^{m}e_a(-1)^a2^{m-i}\alpha^{m+1-a}\\\\
=\sum_{a=1}^me_a\cdot(-1)^a\cdot \alpha^{m+1-a} +\sum_{i=1}^m(-1)^{i+1}2^{m-i}\alpha^{m+1-i}e_i +\sum_{a=1}^{m-1}e_a(-1)^a \alpha^{m+1-a}\sum_{i=a+1}^{m}2^{m-i}.
\end{gather*}
The inner sum over $i$ in the third term is now a geometric series with value $2^{m-a}-1$,
\begin{gather*}
= \sum_{a=1}^me_a\cdot(-1)^a\cdot \alpha^{m+1-a}+\sum_{i=1}^m(-1)^{i+1}2^{m-i}\alpha^{m+1-i}e_i+\sum_{a=1}^{m-1}e_a(-1)^a \alpha^{m+1-a}\cdot\big(2^{m-a}-1\big)
\end{gather*}
\begin{gather*}
= \sum_{a=1}^me_a\cdot(-1)^a\cdot \alpha^{m+1-a}+\sum_{i=1}^m(-1)^{i+1}2^{m-i}\alpha^{m+1-i}e_i\\
+\sum_{a=1}^{m-1}e_a(-1)^a \alpha^{m+1-a}\cdot 2^{m-a}-\sum_{a=1}^{m-1}e_a(-1)^a \alpha^{m+1-a}.
\end{gather*}
The second and third sums are identical with opposite signs except the second sum has an extra $m$ term, so the overall contribution of sum two + sum three is just the $i=m$ term of sum two:
\begin{gather*}
= \sum_{a=1}^me_a\cdot(-1)^a\cdot \alpha^{m+1-a}+(-1)^{m+1}2^0\alpha^{1}e_m -\sum_{a=1}^{m-1}e_{a}(-1)^a \alpha^{m+1-a}.
\end{gather*}
The term that just appeared is the negative of the $a=m$ term of (the new) sum three. Since the sum comes with a - sign, it can then be absorbed 
\begin{gather*}
= \sum_{a=1}^me_{a}\cdot(-1)^a\cdot \alpha^{m+1-a} -\sum_{a=1}^{m}e_{a}(-1)^a \alpha^{m+1-a} \equiv 0.
\end{gather*}
\end{proof}
Returning to the proof of Lemma (\ref{Sm determinant}), we can write 
\begin{gather*}
S_{(m)}=h_m +\cancel{(-)} + \sum_{i=1}^m\left(\left[\sum_{k=1}^{m-i}2^{k-1}\alpha^{k+i}h_{m-i-k}\right]+\alpha^ih_{m-i}\right)\\
=h_m +\sum_{i=1}^m\sum_{k=1}^{m-i}2^{k-1}\alpha^{k+i}h_{m-i-k}+\sum_{i=1}^m\alpha^ih_{m-i}
\end{gather*}
where $\cancel{(-)}$ is the term which is shown to be equal to 0 in Lemma (\ref{sum=0}). If we reindex $j:=k+i$, 
\begin{gather*}
=h_m + \left(\sum_{j=2}^m\sum_{i=1}^{j-1}2^{j-i-1}\alpha^jh_{m-j}\right)+\sum_{i=1}^m\alpha^ih_{m-i}\\
=h_m + \left(\sum_{j=2}^m\alpha^jh_{m-j}\sum_{i=1}^{j-1}2^{j-i-1}\right)+\sum_{i=1}^m\alpha^ih_{m-i}.
\end{gather*}
The inner sum over $i$ is now a geometric series with value $2^{j-1}-1$,
\begin{gather*}
=h_m + \left(\sum_{j=2}^m\alpha^jh_{m-j}(2^{j-1}-1)\right)+\sum_{i=1}^m\alpha^ih_{m-i}\\
=h_m + \sum_{j=2}^m2^{j-1}\alpha^jh_{m-j}-\sum_{j=2}^m\alpha^jh_{m-j}+\sum_{i=1}^m\alpha^ih_{m-i}\\
=h_m + \sum_{j=2}^m2^{j-1}\alpha^jh_{m-j}+\alpha^1h_{m-1}\\
=h_m + \sum_{j=1}^m2^{j-1}\alpha^jh_{m-j}
\end{gather*}
as required.
\end{proof}

\end{document}
